%% kapa: exam=1 round=30 year=2019 period=1 form=A subject=real_estate subject_ko=부동산학원론 seq=25 number=105
%% source: Q-Net 감정평가사 기출문제 (cst00309) artlSeq=5205205 file=extracted/1008/5205205/2019년 제30회 감정평가사 1차시험 문제지/2019년 제30회 감정평가사 1차시험 문제지/2019년도 제30회 감정평가사 1차 1교시 A형.hwp
%% conversion: kordoc 4.24.51 -> tools/gen1.py (auto); review=proofread-sonnet; flags=none
\begin{question}{105}
다음의 조건을 가진 A부동산에 관한 설명으로 옳지 않은 것은? (단, 주어진 조건에 한함)
\begin{framedtext}
○ 가능총소득: 연 1억원 \par
○ 공실 및 대손: 가능총소득의 10 \% \par
○ 운영경비: 유효총소득의 30 \% \par
○ 가격구성비: 토지 40 \%, 건물 60 \% \par
○ 토지환원율: 연 3 \%, 건물환원율: 연 5 \%
\end{framedtext}
\begin{choices}
  \item 유효총소득은 연 9천만원이다.
  \item 순영업소득은 연 6천3백만원이다.
  \item 자본환원율은 연 4\%이다.
  \item 수익가격은 15억원이다.
  \item 운영경비는 연 2천7백만원이다.
\end{choices}
\end{question}
\answer{3}
